\documentclass[conference]{IEEEtran}
\usepackage[T1]{fontenc}
\usepackage{amsmath}

\begin{document}
\title{CPE380 Multi-Cycle Team Project:}
\author{
    \IEEEauthorblockN{Benjamin Allison}
    \IEEEauthorblockA{Computer Science \& Computer Engineering}
    \IEEEauthorblockA{University of Kentucky\\
    Lexington, KY USA\\
    baal242@uky.edu}
    \and
    \IEEEauthorblockN{Mohandass Kannan}
    \IEEEauthorblockA{Computer Science \& Computer Engineering \& Artificial Intelligence}
    \IEEEauthorblockA{University of Kentucky\\
    Lexington, KY USA\\
    mka377@uky.edu}
    \and
    \IEEEauthorblockN{Author 3}
    \IEEEauthorblockA{Computer Science/Engineering}
    \IEEEauthorblockA{University of Kentucky\\
    Lexington, KY\\
    abcd123@uky.edu}
}

\maketitle

\section

\\THIS IS AN EXAMPLE SECTIONXXXXXXXXXXX. 
    It works as normal and just \texttt{} for italics for code. 
\section{The \texttt{ori} and \texttt{min} Instructions} 
We decoded \texttt{ori} by adding \texttt{`DECODE(`OP(-1), `OP(13), 24)} to
\texttt{ONOP}, and it runs in states 24--28. However the \texttt{IRimmedout} 
sign-extends the immediate, but MIPS \texttt{ori} is
supposed to zero-extend it. So we loaded \texttt{0x0000ffff} into \texttt{Y} and
mask the immediate with \texttt{ALUand} before the OR statement.
Without that mask, any immediate with bit 15 set would turn the whole upper 
half of the results into ones.
No temporary register needed for this process. For \texttt{min} we added
\texttt{`DECODE(`OP(-1)+`FUNCT(-1), `FUNCT(1), 29)} and used states 29--39.

\section{Part2}

\section{Part3]
